% ID: MAT-ALG-RAZ-003
% Titolo: Espressione con esponenti negativi
% Disciplina: Matematica
% Area: Algebra
% Argomento: Numeri razionali
% Classe: Prima secondaria di secondo grado
% Difficolta: 3
% Tipo: problema_numerico
% Tag: matematica, algebra, numeri_razionali
% Fonte: esercizi_razionali_anonimizzati
% Provenienza: materiale fornito dal docente; testo anonimizzato; soluzione redatta per Exergo

\begin{esercizio}
Semplifica $\left[6\left(\dfrac34\right)^{-2}\right]^{-1}:\left[\left(\dfrac29\right)^{-2}:\left(\dfrac32\right)^4\right]^{-1}+2^{-3}\dfrac{1}{3^{-3}}3^{-2}$.
\end{esercizio}

\begin{soluzione}
La prima parentesi inversa vale $3/32$, la seconda $1/4$. Il primo quoziente vale $3/8$; anche il secondo addendo è $(1/8)27(1/9)=3/8$. Somma: $3/4$.
\end{soluzione}
